% Standalone problem document, generated from the adjacent README.md.
% Compile with XeLaTeX. Mathematical content is maintained in Markdown.
\documentclass[11pt,a4paper]{article}
\usepackage[fontsize=10.5pt]{scrextend}
\usepackage[margin=22mm,headheight=15pt,headsep=9mm,footskip=11mm]{geometry}
\usepackage{fontspec}
\setmainfont{texgyrepagella}[Extension=.otf,UprightFont=*-regular,BoldFont=*-bold,ItalicFont=*-italic,BoldItalicFont=*-bolditalic]
\setsansfont{texgyreheros}[Extension=.otf,UprightFont=*-regular,BoldFont=*-bold,ItalicFont=*-italic,BoldItalicFont=*-bolditalic]
\setmonofont{texgyrecursor}[Extension=.otf,UprightFont=*-regular,BoldFont=*-bold,ItalicFont=*-italic,BoldItalicFont=*-bolditalic,Scale=MatchLowercase]
\usepackage{amsmath,amssymb}
\usepackage{unicode-math}
\setmathfont{texgyrepagella-math.otf}
\usepackage{xcolor,microtype,enumitem,needspace,fancyhdr}
\usepackage{bookmark}
\definecolor{ink}{HTML}{17344B}
\definecolor{accent}{HTML}{197D80}
\definecolor{muted}{HTML}{596773}
\hypersetup{colorlinks=true,linkcolor=accent,urlcolor=accent,pdftitle={TR-15 - Nonnegative
H-eigenvalue inheritance from odd-order Hankel
tensors},pdfauthor={Open Problems in Numerical Linear Algebra}}
\urlstyle{same}
\setlength{\parindent}{0pt}
\setlength{\parskip}{5pt plus 1pt minus 1pt}
\linespread{1.04}
\setlength{\emergencystretch}{3em}
\widowpenalty=10000
\clubpenalty=10000
\displaywidowpenalty=10000
\allowdisplaybreaks[1]
\setlist{nosep,leftmargin=1.5em}
\providecommand{\tightlist}{\setlength{\itemsep}{0pt}\setlength{\parskip}{0pt}}
\providecommand{\pandocbounded}[1]{#1}
\pagestyle{fancy}
\fancyhf{}
\fancyhead[L]{\sffamily\footnotesize\color{muted}OPEN PROBLEMS IN NUMERICAL LINEAR ALGEBRA}
\fancyhead[R]{\sffamily\bfseries\footnotesize\color{accent}TR-15}
\fancyfoot[L]{\parbox[b]{0.9\textwidth}{\sffamily\fontsize{8}{10}\selectfont\color{muted}Editorial ratings; a bounded search cannot certify that a problem remains unsolved.\\Literature check: 2026-09-12}}
\fancyfoot[R]{\sffamily\footnotesize\color{muted}\thepage}
\renewcommand{\headrulewidth}{0.3pt}
\renewcommand{\headrule}{\hbox to\headwidth{\color{accent}\leaders\hrule height \headrulewidth\hfill}}
\makeatletter
\renewcommand{\section}{\Needspace{5\baselineskip}\@startsection{section}{1}{0pt}{12pt plus 2pt minus 2pt}{4pt}{\sffamily\bfseries\large\color{ink}}}
\renewcommand{\subsection}{\Needspace{4\baselineskip}\@startsection{subsection}{2}{0pt}{9pt plus 1pt minus 1pt}{3pt}{\sffamily\bfseries\normalsize\color{ink}}}
\makeatother
\setcounter{secnumdepth}{0}
\begin{document}
{\sffamily\small\color{muted}Tensor computations\par}
\vspace{3pt}
{\raggedright\hyphenpenalty=10000\exhyphenpenalty=10000\sffamily\bfseries\fontsize{21}{24}\selectfont\color{ink}Nonnegative
H-eigenvalue inheritance from odd-order Hankel tensors\par}
\vspace{7pt}
{\sffamily\small\textbf{Difficulty:} challenging\quad\textbf{Importance:} interesting
to specialist\par}
{\sffamily\small\bfseries\color{accent}Status: Lean verified\par}
\vspace{4pt}
{\color{accent}\hrule height 0.8pt}
\vspace{6pt}
\textbf{Rating rationale:} Challenging because odd-order H-eigenvalue
positivity lacks the even-order polynomial-positivity argument;
specialist importance lies in transferring spectral certificates within
Hankel representations.

\subsection{Lean proof and verification evidence -
2026-09-12}\label{lean-proof-and-verification-evidence---2026-09-12}

\textbf{The complete odd-order inheritance conjecture is false, with a
Lean-verified counterexample.} The
\href{https://github.com/sgstepaniants/OpenProblemsInNLA/tree/6a2d0868e8b7905dbd5c04d4beaae8cf43288e1f/tensor-computations/TR-15/lean}{proof
at revision 6a2d086} uses \(m=3\), \(q=2\), \(n=2\) and the shared
generator \(h=(2,0,1,0,2,0,-1)\). Every real H-eigenvalue of the lower
tensor is strictly positive, while the upper tensor has the exact
H-eigenpair \((-1,(0,1))\).

\textbf{Mathematical counterexample and informal proof:} Matthew J.
Colbrook, Department of Applied Mathematics and Theoretical Physics,
University of Cambridge. \textbf{Lean formalization:} George
Stepaniants, Department of Computing and Mathematical Sciences,
California Institute of Technology, Pasadena, California, USA, with
AI-agent assistance.

The definitions retain every original odd \(m\ge3\), integer \(q\ge2\),
dimension \(n\ge2\) and real common generator. They use the actual
Hankel entries, every ordered contraction tuple, signed coordinate
powers and nonzero real H-eigenvectors. The seven
\href{https://github.com/sgstepaniants/OpenProblemsInNLA/blob/6a2d0868e8b7905dbd5c04d4beaae8cf43288e1f/tensor-computations/TR-15/lean/Solution.lean}{checked
exports}, each with prefix \textit{NLA.TR15.}, are:

\begin{itemize}
\tightlist
\item
  \textit{lower\_contractions}: all three exact component polynomials
  for every real vector.
\item
  \textit{upper\_contraction}: the actual order-six contraction at
  \((0,1)\).
\item
  \textit{lower\_eigenvalues\_pos}: positivity of every lower real
  H-eigenvalue.
\item
  \textit{lower\_eigenpair\_exists}: a genuine nonzero lower pair from
  the intermediate value theorem.
\item
  \textit{upper\_negative\_eigenpair}: the complete upper pair and
  strict negative eigenvalue.
\item
  \textit{counterexample}: original admissibility, the lower premise and
  failure of the upper conclusion.
\item
  \textit{not\_inheritanceConjecture}: negation of the complete
  universal inheritance assertion.
\end{itemize}

Lower positivity follows from the first contraction, a strictly positive
sum of squares. The upper contraction has one surviving ordered tuple.
An actual intermediate-value root supplies the lower pair; its existence
is proved without adding a premise or a strong-Hankel hypothesis.

Two independent agents approved the
\href{https://github.com/ajt60gaibb/OpenProblemsInNLA/blob/main/tensor-computations/TR-15/lean/reviews}{statements
and completed proof}.
\href{https://github.com/sgstepaniants/OpenProblemsInNLA/actions/runs/34716902324}{Linux
run 34716902324} matched all seven exports with the sandboxed Comparator
and replayed the solution in Lean's default kernel. The
\href{https://github.com/ajt60gaibb/OpenProblemsInNLA/blob/main/tensor-computations/TR-15/lean/verification/linux-2026-09-12}{original
artifacts and independent audit} bind all 135 input hashes and both
actual isolation/rejection-control suites. All 15
\href{https://github.com/ajt60gaibb/OpenProblemsInNLA/blob/main/tensor-computations/TR-15/lean/verification/linux-2026-09-12/axiom-verification.json}{internal/public
transitive axiom reports} use only \textit{propext},
\textit{Classical.choice} and \textit{Quot.sound}. These are independent
agent reviews, without a claim of external human peer review.

The proof pins \textbf{Lean 4.33.1},
\href{https://github.com/alerad/leancert/tree/621a43d7cf21f87872392a01e874f2f1dbddc926}{LeanCert
621a43d} and
\href{https://github.com/leanprover-community/mathlib4/tree/0df444a360eaa60ab8c11dca51a86af692955474}{Mathlib
0df444a}. The kernel-mode LeanCert certificate \(-1<0\) remains in the
final negative-pair contradiction. See the
\href{https://github.com/ajt60gaibb/OpenProblemsInNLA/blob/main/tensor-computations/TR-15/lean/formalization.yaml}{manifest},
\href{https://github.com/ajt60gaibb/OpenProblemsInNLA/blob/main/tensor-computations/TR-15/lean/lake-manifest.json}{dependency
pins} and
\href{https://github.com/ajt60gaibb/OpenProblemsInNLA/blob/main/tensor-computations/TR-15/lean/NUMERICAL_TARGETS.md}{exact
targets}. From the verified revision on a documented
\href{https://github.com/ajt60gaibb/OpenProblemsInNLA/blob/main/tools/lean/HARNESS.md}{non-root
Linux host}, reproduce with:

\begin{verbatim}
tools/lean/bootstrap.sh /absolute/path/to/nla-lean-tools
tools/lean/selftest.sh /absolute/path/to/nla-lean-tools
tools/lean/verify.sh \
  tensor-computations/TR-15/lean \
  /absolute/path/to/nla-lean-tools
\end{verbatim}

\subsection{Resolution --- 2026-09-11}\label{resolution-2026-09-11}

\textbf{Negative resolution.} Matthew J. Colbrook's
\href{https://github.com/ajt60gaibb/OpenProblemsInNLA/blob/main/references/colbrook-unclaimed-2026-09-11/manuscripts/TR-15.md}{complete
manuscript, Counterexample proposition and proof} disproves the
universal conjecture with \(m=3\), \(q=2\), \(n=2\) and the common
generating vector \(h=(2,0,1,0,2,0,-1)\). Every real H-eigenvalue of the
order-three, dimension-three tensor \(A\) is strictly positive, whereas
the order-six, dimension-two tensor \(B\) has the exact H-eigenpair
\((-1,(0,1))\). The manuscript also proves that \(A\) has a real
H-eigenpair, so its premise is nonvacuous. The counterexample meets the
original odd-lower-order assumptions; it does not contradict the
separate even-lower-order theorem or results requiring a
positive-semidefinite associated Hankel matrix.

The complete source passed
\href{https://github.com/ajt60gaibb/OpenProblemsInNLA/blob/main/references/colbrook-unclaimed-2026-09-11/verification/reviews/TR-15-review.md}{independent
Codex-agent proof review}.
\href{https://github.com/ajt60gaibb/OpenProblemsInNLA/blob/main/references/colbrook-unclaimed-2026-09-11/README.md}{Authorship,
AI assistance and verification record}. The 2026-09-11 review did not
assert external human peer review or formal verification. The later Lean
verification above covers the complete original implication and proves
the premise is nonvacuous. The original target below is retained
verbatim; the difficulty, importance and rating rationale above are
historical. This entry no longer contributes to the open count.

\href{https://github.com/ajt60gaibb/OpenProblemsInNLA/blob/main/references/colbrook-unclaimed-2026-09-11/manuscripts/TR-15.pdf}{Manuscript
PDF}.

\subsection{Statement}\label{statement}

For every odd \(m\ge3\), integer \(q\ge2\), dimension \(n\ge2\), and
vector \(h\in\mathbb R^{qm(n-1)+1}\), define Hankel tensors \(A\) and
\(B\) with this same generating vector. Their orders and dimensions are

\[\mathop{\mathrm{order}}\nolimits(A)=m,\quad\dim(A)=q(n-1)+1,
\qquad \mathop{\mathrm{order}}\nolimits(B)=qm,\quad\dim(B)=n.\]

An order-\(s\), dimension-\(N\) Hankel tensor has entry
\(h_{i_1+\cdots+i_s-s}\), with each index in \(\{1,\ldots,N\}\).

For a real order-\(s\) tensor \(C\), an H-eigenvalue is a real
\(\lambda\) admitting a real nonzero vector \(x\) such that, for every
\(i\),

\[\sum_{i_2,\ldots,i_s=1}^N C_{i i_2\cdots i_s}x_{i_2}\cdots x_{i_s}
=\lambda x_i^{s-1}.\]

Conjecture: if \(A\) has no negative H-eigenvalues, then \(B\) has no
negative H-eigenvalues.

\subsection{Relevance}\label{relevance}

The question transfers a spectral positivity certificate between
different Hankel representations of the same data, relevant to
structured tensor eigenvalue computation and polynomial positivity.

\subsection{References}\label{references}

\begin{enumerate}
\def\labelenumi{\arabic{enumi}.}
\tightlist
\item
  W. Ding, L. Qi, and Y. Wei, \emph{Inheritance properties and
  sum-of-squares decomposition of Hankel tensors: theory and
  algorithms}, BIT Numer. Math. 57 (2017), 169--190.
  \href{https://doi.org/10.1007/s10543-016-0622-0}{DOI};
  \href{https://www.polyu.edu.hk/ama/staff/new/qilq/BIT-DQW.pdf}{author
  PDF}, final §4, ``The third inheritance property of Hankel tensors,''
  concluding conjecture; §2 gives order/dimension conventions.
\item
  L. Qi, \emph{Hankel Tensors: Associated Hankel Matrices and
  Vandermonde Decomposition}, 2014.
  \href{https://arxiv.org/pdf/1310.5470}{Primary preprint}, §5,
  concerning H-eigenvalues and complete Hankel tensors.
\end{enumerate}

\subsection{Status check --- 2026-09-10}\label{status-check-2026-09-10}

Rechecked the
\href{https://www.polyu.edu.hk/ama/staff/new/qilq/BIT-DQW.pdf}{Ding--Qi--Wei
journal text, final §4}, and searched for proofs or counterexamples to
the third inheritance property. The source explicitly leaves odd lower
order unresolved; its even-lower-order result lies outside the displayed
target. Complete or strong Hankel hypotheses would be additional
restrictions. No later resolution was located; this remains a bounded
historical-source check.
\end{document}
